\section{Related work}\label{sec:related}
\paragraph{SQ learning and dimension.}
Statistical-query learning accesses bounded expectations rather than individual examples \citep{Kearns98}. Feldman's randomized search characterization supplies a general dimension of this access model \citep{Feldman17}; our use of it is confined to Appendix~\ref{sec:rsd}. The question of \citet{FKS26} asks whether one learner serving every marginal must admit a common low-dimensional linear representation. Our rectangle learner answers negatively without invoking the characterization. Euclidean representation complexity originates in \citet{BDES02}; probabilistic and approximate versions were developed by \citet{CMW25} and \citet{KMS20}. The exact and averaged variants require the distinct reductions in \cref{thm:prob,lem:exact-prob}. In particular, uniform approximation on the sparse random matrix would not prove the averaged-dimension lower bound.

\paragraph{Correlational queries and the converse direction.}
The source note records a margin-complexity characterization of distribution-independent weak correlational learning \citep[Section~3]{FKS26}, citing \citet{Feldman08}. Strong correlational learning is more restrictive: \citet{Feldman11} proves hardness for conjunctions at vanishing target error. These two accuracy regimes are distinct. General SQ halfspace algorithms use rescaling and marginal information \citep{DV04,BF13}. The rescaling simulation depends polynomially on dimension, inverse accuracy and the logarithm of the inverse margin; finite-grid margin bounds control the latter through coordinate bit complexity \citep{GSS13}. The negative result here concerns reversing the implication without requiring a low-dimensional or bounded-margin representation.

\paragraph{Prior random-feature bounds.}
\citet{KM25} study approximate dimension complexity, random features and distribution-free differentiable learning. The source note states that their claimed upper bounds do not hold, attributing the correction to a 2026 personal communication from those authors \citep[Section~3]{FKS26}. The probabilistic statement in \cref{cor:negative-answer} fixes one embedding law before both the target and marginal. Under these quantifiers it rules out an unrestricted polynomial bound from common SQ learnability to the nondegenerate averaged dimension. This comparison is about the precisely stated implication, rather than an alternative marginal-dependent random-feature distribution.

\paragraph{Sign-rank, margin and geometry.}
Forster's normalization method provides spectral lower bounds \citep{Forster02}. The construction and the separation between sign-rank, VC dimension, average margin and rectangle ratio are due to \citet{HHPTZ22}, building on \citet{AMY16}. We import their grid construction and Warren's general strict-pattern theorem \citep{Warren68}; we derive the application-specific counting and probabilistic lower bounds. The homogeneous-rectangle theorem of \citet{APP05} is used only to sharpen the inverse rectangle constant to $2^{15}$. Our cap proof supplies an independent dimension-dependent constant. The minimax rectangle-mixture idea occurs in the average-margin argument of \citet[Theorem~3.1]{HHPTZ22}; the new step here is its adversarial-oracle peeling/elimination analysis and bounded query budget. Approximate Forster transforms are prior algorithmic work \citep{DKT22}; Appendix~\ref{app:convex} gives the fixed-dimensional implementation needed here.

\paragraph{The finite-field separation of 2017.}
\citet[Theorem~7.7]{Feldman17} studies affine lines on $\mathbb F_p^2$ and proves a distribution-independent weak-learning lower bound of $t/2-1$ queries at tolerance $1/t$, where $t=(p/32)^{1/4}$, for the stated accuracy and success-probability regime. Theorem~7.8 there gives a separately fixed-distribution learner with $O(1/\epsilon^2)$ queries of tolerance $\epsilon^2/13$. For the affine matrix, Proposition~\ref{prop:affine} applies the joint certificate to his affine class. At constant expected error it gives $m/\tau^2=\Omega(p)$, hence $m=\Omega(p)$ at constant tolerance. The proposition also supplies the exact success-probability comparison at his tolerance. These are quantitative refinements of the hardness side; his class is not a small-complexity distribution-independent learner and therefore is not the present counterexample. Our real-grid monotone flips lie on the easy side of that universal-marginal requirement.

\paragraph{Approximate sign-rank, May 2026.}
\citet[Theorem~1.3]{BHKR26} extend the large-rectangle property to approximate sign-rank and partial matrices. Their approximate rank fixes one \emph{deterministic} column embedding and, for each row and each marginal, permits a new separating vector with error at most $\eta$. It is not the random-embedding quantity $\dc^\eta$ used here. Their theorem supplies row and column fractions $d^{-C_\eta d}$ and $d^{-C_\eta d^2}$, respectively, with $C_\eta$ stated as $O(\log(1/(1-2\eta)))$, for fixed $0<\eta<1/2$; one may enlarge the constant and use $d\ge2$. \Cref{cor:approxrect} converts this published geometric theorem into a common SQ learner with the displayed parameters. For partial matrices their rectangle is homogeneous on nonstar entries, possibly containing stars; we do not silently regard it as a full monochromatic rectangle of an arbitrary completion.

\paragraph{Index and list replicability, June 2026.}
\citet{BHHLT26} prove $\operatorname{Ind}_{\mathbb Z_2}(H)\le2\operatorname{LR}(H)-1$ and show that projective-plane incidence classes have list replicability at most three and index at most five despite unbounded sign-rank. Our joint-distribution theorem shows that those same projective-plane classes are SQ-hard when a common learner must serve every marginal. Thus bounded list replicability or bounded index is not the missing SQ upper-bound criterion. Their separation concerns different parameters and does not resolve the question answered here. Earlier and concurrent topological bounds \citep{HHM23,FHV26}, and the total Hamming-distance upper bounds of \citet{GHIS25}, likewise do not give the sought efficient explicit high-sign-rank selection for the constant-rectangle flip family.

\paragraph{The August 2026 conditional SQ theorem.}
A further comparison is \citet[Assumption~4.42 and Theorem~4.9]{VALG26}. Its SQ result assumes that the span of all seed-averaged terminal response functions, over complete response rules, already has dimension at most $B(1+m/\tau^2)^k$. Under that additional assumption it proves a dimension upper bound. The extra response-rank hypothesis is not a consequence of the SQ interface proved there. Our counterexample rules out deriving such a universal polynomial hypothesis from distribution-independent SQ learnability alone. For the deterministic learners in \cref{thm:caps}, the universal mean-response-rank hypothesis would imply a polynomial dimension upper bound and hence contradict \cref{thm:construction}. Thus it cannot follow from the common SQ guarantee alone. 

\paragraph{Learning versus communication and differentiable learning.}
The converse communication construction in Theorem~\ref{thm:communication} is \citet[Proposition~3.14]{HHPTZ22}; its composition with RECTIFY yields the SQ sufficient condition and the halfspace boundary case. The proper completion and the order-median succinct learner exploit the real incidence geometry rather than invoking a generic protocol-to-learning conversion. The differentiable-learning simulations of \citet{AKMSS21} use engineered models and initialization and charge the simulated computation in their network-size bounds. They do not preserve the remaining prescribed benign-SGD algorithm. The hidden-key bounds in Section~\ref{sec:secret} explicitly distinguish uniform key-independent learners from class-dependent architecture choices.


\paragraph{Neural dynamics and staircase structure.}
\citet{MKAS21} compare gradient learning with tangent kernels under specified model conditions. The staircase analyses of \citet{ABBBN21} and \citet{ABM22} identify hierarchical structures learnable by particular regular-network dynamics under their distributional assumptions. Their conclusions concern those dynamics and data regimes. \Cref{cor:engineered} instead composes a common SQ learner with the full-batch simulation of \citet{AKMSS21}; its engineered architecture and initialization retain the simulation's computational cost. The prescribed Gaussian-initialized online-SGD question therefore remains a separate problem.
